\subsection{Local Expansion in Low-Density Environments}
  \label{subsec:local_expansion}

  Measurements of the Hubble constant rely on local kinematic observables,
  including distance ladders, standard candles, and redshift--distance relations.
  These measurements implicitly assume that the locally inferred expansion rate
  is representative of the global cosmological background~\cite{Riess2022SH0ES,Freedman2021TRGB}.

  However, observations of large-scale structure show that the late-time Universe
  is strongly inhomogeneous.
  A substantial fraction of the cosmic volume is occupied by voids,
  characterized by matter densities well below the cosmic mean~\cite{Keenan2013KBC}.

  Within the effective framework introduced in
  Section~\ref{sec:effective_model},
  low-density environments probe the saturated regime of the effective response.
  In this regime, local kinematic observables no longer sample the same effective
  transition resolution as in high-density regions.
  As a result, expansion rates inferred from such environments need not coincide
  with the global average, even in the absence of any modification of early-time
  cosmology.

  We therefore distinguish between a global expansion rate
  $H_{\mathrm{global}}$,
  defined by the homogeneous background constrained by early-universe
  observables~\cite{Planck2020},
  and an effective local expansion rate $H_{\mathrm{loc}}$,
  inferred from observations within a specific late-time environment.
